\chapter{Supply-Chain Integrity and Reproducible Builds}
\label{ch:supply}

\section{Motivation}
\label{sec:supply-mot}

Verification is worthless if the artifact a user downloads was not the one
that was verified. This chapter addresses the supply chain: how MATHLIB5
makes its build reproducible and how the anchor (\cref{ch:anchors}) lets a
third party confirm they received the verified artifact.

\section{Reproducibility Properties}
\label{sec:supply-repro}

The build is designed to be bit-for-bit reproducible given the same
inputs:
\begin{itemize}
  \item Source is pinned via the Git commit that the WORM receipt
        references.
  \item Toolchain versions are recorded in \texttt{package.json} and the
        lockfile, so \texttt{npm ci} yields identical dependencies.
  \item Deterministic ordering of proof search avoids nondeterministic
        output in the receipts.
\end{itemize}
A rebuild on a different machine that produces a different artifact hash is
therefore a signal of tampering or environment drift, not expected
behavior.

\section{The Anchor as a Trust Transfer}
\label{sec:supply-anchor}

Because the PDF and the repository both hash into the same anchor chain,
trust in the public Bitcoin transaction transfers to the local artifact:
\begin{enumerate}
  \item The verifier fetches the Bitcoin \texttt{OP\_RETURN} and recovers
        the expected SHA-256.
  \item The verifier hashes the local artifact and compares.
  \item Equality implies the local artifact is the one that was publicly
        anchored, with no trusted intermediary required.
\end{enumerate}

\section{Dependency Attestation}
\label{sec:supply-dep}

Each external dependency used during the build is itself hashed and
recorded in the receipt metadata, so a compromised dependency (a classic
supply-chain vector) changes the overall artifact hash and is caught by
the anchor comparison. This is the practical content of the
AGENTS.md claim that the repo is the substrate: the repository's hash
\emph{is} the unit of trust.

\section{Limits}
\label{sec:supply-limits}

Reproducibility protects the \emph{software} artifact, not the
\emph{hardware} it runs on; a malicious CPU can still misexecute verified
C$--$. The anchor also assumes the Bitcoin ledger's integrity
(\cref{ch:threat}). Within these standard assumptions, however, the chain
gives end users a way to confirm they hold exactly the verified system.
